% Study 5: supplementary methods, the sample, and the results that no other part of the paper
% prints. Numbers come only from macros in generated/numbers.tex.
\label{sec:suppB-study5}

Study~5 asked whether the hours Alex was said to have spent making a second work change its ratings, or change the AI label's effect on them. Its data were collected from 21 October to 2 November 2025, after the CogSci 2025 paper and the September 2025 talk, neither of which reports it. Its registration, a dated record of its hypotheses and analysis plan, was filed on the Open Science Framework (OSF) on 8 October 2026. That was after data collection but before any rating had been averaged or compared by label (Supplementary Table~\ref{tab:supp-verdicts}). The main text reports the registered tests; Supplementary Note~5 and Supplementary Fig.~\ref{figS:editingtime} break down the one confirmatory test that was significant after Holm's correction; and Supplementary Table~\ref{tab:supp-labeleffects} gives the label's effect at each level. This section gives the materials, procedure and sample, and the results that no other part of the paper prints.

\subsection*{Materials}

Study~5 used the 30 images of Studies 2 and 4: ten originals, all generated with OpenAI's DALL·E, an AI image generator, each with a Slight and a Dramatic version that we made by hand and a None second work, the original shown again unchanged (Section~\ref{sec:suppA-registrations}; Supplementary Note~4). The Human label was therefore a cover story. Each participant saw six of the ten originals, three with a human or humanoid figure and three without, each followed by its three second works.

The labels and the vignette before each original were worded as in Study~4. The Human label read: `There is a digital artist named Alex. Alex is a skilled artist who creates digital artworks. Alex uses his graphic design skills and digital sketching tools.' The AI label read: `There is a digital artist named Alex. Alex is a skilled artist who creates digital artworks. Alex develops and inputs prompts into an artificial intelligence (AI) system. The AI system generates digital artworks based on those prompts.' Before each original, the Human group read `Alex had a new idea for a digital artwork. Alex spent \textit{time} working with his graphic design and digital sketching tools, revising as his work as he goes, which creates a digital artwork.' The AI group read `Alex had a new idea for a work of digital artwork. Alex spent \textit{time} developing and inputting prompts into the AI image generator, revising each prompt as he goes, which creates a digital artwork.' Here \textit{time} was `less than 1 hour', `10 hours' or `100 hours', two originals at each. The survey showed this time on every trial but did not record it, because the field that should have stored it had been renamed in the survey. It therefore enters no analysis.

The description of each second work now stated the hours Alex spent making it. Before the first of the three, both groups read `Alex looked at the first artwork. He decided to produced a secondary digital artwork. Using the first artwork as a base, he spent \textit{editing time} producing a second digital artwork with his own graphic design skills and digital sketching tools.' Before each of the other two, they read `Later, Alex looked at the very first artwork again. Alex decided to produce a different secondary digital artwork. Using the very first artwork as a base, he spent \textit{editing time} producing a different secondary digital artwork with his own graphic design skills and digital sketching tools.' Each ended `Click the arrow to see the digital artworks:'. The editing time was `less than 1 hour', `10 hours' or `100 hours', and it was the same for a given participant throughout. Because the sentence was the same at every level, at None it said that Alex had spent the stated hours on an unchanged copy (Supplementary Note~5). The quotations reproduce the survey exactly, slips included.

The three ratings, their definitions, their answer scales, the layout of the rating screen and the required explanations were as in Study~4, and a script again blocked copying and pasting in 21 of the survey's 90 explanation boxes and not in the others. Two transformativeness items, both for one original, had six answer points. As registered, their \SFiveNSixPointDropped{} ratings were left out of the primary models and entered only the rescaled analysis described below. The slips in three essence answer labels and one creativity question recurred; scoring by answer position is unaffected.

Study~5 thus differs from Study~4 in two ways. A stated editing time was added, and it varied between participants, whereas Study~4's stated time for the original varied within them. And the stated time for the original, though shown, was not recorded.

\subsection*{Procedure}

Participants took part online in a Qualtrics survey. By our account, participants were students in the UCLA Psychology subject pool who took part for course credit. Every analyzed record carries the code that the survey sent, at its end, to the pool's page for granting credit. The consent form, the same in all five studies, also offered hourly pay to people who signed up outside the pool, and no record shows that anyone took that route. The survey assigned each participant at random to one of the six combinations of label and editing time for the whole session, presenting each combination equally often. It also set a quota of 40 participants for each combination, and it never moved anyone from one combination to another. Each original was shown alone for at least 8 seconds, and each second work beside it for at least 10 seconds. The order of the three levels was random and was not recorded. Each participant rated 18 second works, giving 54 ratings and 54 explanations.

The task question, the arithmetic question (after the trials) and the closing screen, which asked the seriousness question before the background questions, were as in Studies 2 to 4. No question asked participants to recall the label (a manipulation check). The debriefing was the same as in the other four studies; Section~\ref{sec:suppA-study1} quotes it and corrects its statement of the law. Participants read it only after their last rating.

\subsection*{Exclusions and sample}

The registered exclusion rules, applied in order, left \SFiveN{} of the \SFiveNRecords{} records in the survey export. There were no survey test runs. Then \SFiveNUnfinishedRemoved{} unfinished records were removed, \SFiveNNotSeriousRemoved{} participants who gave any answer but `Yes, I did' to the seriousness question, a blank included, \SFiveNArithmeticRemoved{} who answered the arithmetic question wrongly, and \SFiveNNoConditionRemoved{} records with no recorded condition, which held no ratings. Of the \SFiveN{} participants analyzed, \SFiveNHuman{} saw the Human label and \SFiveNAI{} the AI label; Supplementary Note~5 gives the six group sizes. They gave \SFiveNRatingsAll{} ratings, each with its explanation, of which \SFiveNRatings{} entered the primary models once the six-point items were left out.

The participants' mean age was \SFiveAgeMean{} years (standard deviation \SFiveAgeSD; range \SFiveAgeMin{} to \SFiveAgeMax). Of the \SFiveN, \SFiveNWomen{} described their gender as female, \SFiveNMen{} as male, and \SFiveNGenderOther{} chose `Other' or `Prefer not to say' or gave no answer. The answers on race, native language and formal artistic training are not summarized here.

Study~5 had its own listing in the subject pool, separate from the one that Studies 1 to 4 shared, and some of its participants had taken part in one earlier study. Across the five studies, the \NTotal{} analyzed records come from \NPeople{} people. Of these, \NPeopleRepeat{} took part in two studies, in each case Study~5 and one earlier study, so they had already seen some of the same originals, and read the same debriefing, before. No registered analysis removes them, and no analysis checks whether their first session changed their ratings. The pooled analyses of Supplementary Notes 6 and 7 and Section~\ref{sec:suppB-crossstudy} count a person who took part in Study~5 and in Study~2 or~4 as two participants.

\subsection*{Results not reported elsewhere}

Each result below carries its status. A registered confirmatory test is one that the registration names as confirmatory, judged by its rule after Holm's correction, which raises each $p$-value in a set of tests so that the chance of calling any of them significant by mistake stays at .05. A $p$-value is the chance of a result at least this large if there were no effect, and a result is called significant when $p$ is below .05. A registered secondary analysis is one the registration names outside that set. An exploratory analysis was named as exploratory or not named at all, and carries no verdict.

\paragraph*{Registered tests and predicted means.}
Supplementary Table~\ref{tab:suppB-s5} gives all six registered confirmatory tests, the registered secondary equivalence tests of editing time on all three measures, and the predicted means. An $F$~test asks whether a factor changes the mean rating. An equivalence test asks whether an effect is shown to lie inside a stated bound; a pass means that the effect is small, not that it is zero. Two registered checks repeated the six confirmatory tests. The first was a cumulative-link mixed model, which treats the five answers as ordered categories rather than evenly spaced numbers; it reached the same verdict as the linear model on all six (Supplementary Note~5 gives its creativity interaction). The second rescaled every item from 0 to 1, so that the two six-point items could be included, and it too gave the same six verdicts. All registered linear, ordinal and rescaled models converged, meaning that the fitting procedure settled on a solution, and none was singular, that is, none estimated a source of variation at zero.

% Study 5: the six registered confirmatory tests with the rescaled analysis, the registered
% secondary editing-time contrasts, and the predicted means.
% Every number is a macro from generated/numbers.tex; supplement/tables/README_suppB.md lists
% the pipeline tables behind them.
\begin{table}[htbp]
\centering
\footnotesize
\caption{\textbf{Study~5: the registered tests of stated editing time, and the predicted means.} Ratings are on 1-to-5 scales. \textbf{a},~The six registered confirmatory $F$~tests, corrected together by Holm's method, beside the registered second analysis with every item rescaled from 0 to 1. \textbf{b},~The registered secondary equivalence tests of editing time, each uncorrected: 100 hours minus under 1 hour, averaged over label and level, against a bound of $\pm$\SFiveEditTimeBound{} points, the size of Study~3's creativity time effect. \textbf{c},~Predicted means (descriptive), each averaged over the factors not named.}\label{tab:suppB-s5}
\setlength{\tabcolsep}{4pt}
\begin{tabular}{@{}lccccc@{}}
\multicolumn{6}{@{}l}{\textbf{a}\enspace Registered confirmatory tests}\\
\toprule
 & & & & \multicolumn{2}{c}{Rescaled analysis}\\
\cmidrule(l){5-6}
Term & $F$ ($df$) & Uncorrected & After Holm & $F$ & After Holm\\
\midrule
\multicolumn{6}{@{}l}{\textit{Transformativeness}}\\
Editing time & \SFiveTransEditTimeF{} (\SFiveTransEditTimeDfNum, \SFiveTransEditTimeDfDen) & $p$~\SFiveTransEditTimeP & $p$~\SFiveTransEditTimePHolm & \SFiveTransEditTimeRescaledF & $p$~\SFiveTransEditTimeRescaledPHolm\\
Label by editing time & \SFiveTransSourceByEditTimeF{} (\SFiveTransSourceByEditTimeDfNum, \SFiveTransSourceByEditTimeDfDen) & $p$~\SFiveTransSourceByEditTimeP & $p$~\SFiveTransSourceByEditTimePHolm & \SFiveTransSourceByEditTimeRescaledF & $p$~\SFiveTransSourceByEditTimeRescaledPHolm\\
\addlinespace[3pt]
\multicolumn{6}{@{}l}{\textit{Essence change}}\\
Editing time & \SFiveEssEditTimeF{} (\SFiveEssEditTimeDfNum, \SFiveEssEditTimeDfDen) & $p$~\SFiveEssEditTimeP & $p$~\SFiveEssEditTimePHolm & \SFiveEssEditTimeRescaledF & $p$~\SFiveEssEditTimeRescaledPHolm\\
Label by editing time & \SFiveEssSourceByEditTimeF{} (\SFiveEssSourceByEditTimeDfNum, \SFiveEssSourceByEditTimeDfDen) & $p$~\SFiveEssSourceByEditTimeP & $p$~\SFiveEssSourceByEditTimePHolm & \SFiveEssSourceByEditTimeRescaledF & $p$~\SFiveEssSourceByEditTimeRescaledPHolm\\
\addlinespace[3pt]
\multicolumn{6}{@{}l}{\textit{Creativity}}\\
Editing time & \SFiveCreatEditTimeF{} (\SFiveCreatEditTimeDfNum, \SFiveCreatEditTimeDfDen) & $p$~\SFiveCreatEditTimeP & $p$~\SFiveCreatEditTimePHolm & \SFiveCreatEditTimeRescaledF & $p$~\SFiveCreatEditTimeRescaledPHolm\\
Label by editing time & \SFiveCreatSourceByEditTimeF{} (\SFiveCreatSourceByEditTimeDfNum, \SFiveCreatSourceByEditTimeDfDen) & $p$~\SFiveCreatSourceByEditTimeP & $p$~\SFiveCreatSourceByEditTimePHolm & \SFiveCreatSourceByEditTimeRescaledF & $p$~\SFiveCreatSourceByEditTimeRescaledPHolm\\
\bottomrule
\end{tabular}

\medskip
\begin{tabular}{@{}lcccc@{}}
\multicolumn{5}{@{}l}{\textbf{b}\enspace Registered secondary tests: 100 hours minus under 1 hour of stated editing time}\\
\toprule
Measure & Estimate & 90\% CI & Two-sided test & Equivalence test\\
\midrule
Transformativeness & \SFiveTransEditTimeContrastEst & \SFiveTransEditTimeContrastLoNinety{} to \SFiveTransEditTimeContrastHiNinety & $p$~\SFiveTransEditTimeContrastP & $p$~\SFiveTransEditTimeContrastPTost\\
Essence change & \SFiveEssEditTimeContrastEst & \SFiveEssEditTimeContrastLoNinety{} to \SFiveEssEditTimeContrastHiNinety & $p$~\SFiveEssEditTimeContrastP & $p$~\SFiveEssEditTimeContrastPTost\\
Creativity & \SFiveCreatEditTimeContrastEst & \SFiveCreatEditTimeContrastLoNinety{} to \SFiveCreatEditTimeContrastHiNinety & $p$~\SFiveCreatEditTimeContrastP & $p$~\SFiveCreatEditTimeContrastPTost\\
\bottomrule
\end{tabular}

\medskip
\begin{tabular}{@{}lccccccccc@{}}
\multicolumn{10}{@{}l}{\textbf{c}\enspace Predicted means}\\
\toprule
 & \multicolumn{2}{c}{None} & \multicolumn{2}{c}{Slight} & \multicolumn{2}{c}{Dramatic} & \multicolumn{3}{c}{Stated editing time}\\
\cmidrule(lr){2-3}\cmidrule(lr){4-5}\cmidrule(lr){6-7}\cmidrule(l){8-10}
Measure & Human & AI & Human & AI & Human & AI & Under 1 hour & 10 hours & 100 hours\\
\midrule
Transformativeness & \SFiveTransMeanHumanNone & \SFiveTransMeanAINone & \SFiveTransMeanHumanSlight & \SFiveTransMeanAISlight & \SFiveTransMeanHumanDramatic & \SFiveTransMeanAIDramatic & \SFiveTransMeanUnderOneHour & \SFiveTransMeanTenHours & \SFiveTransMeanHundredHours\\
Essence change & \SFiveEssMeanHumanNone & \SFiveEssMeanAINone & \SFiveEssMeanHumanSlight & \SFiveEssMeanAISlight & \SFiveEssMeanHumanDramatic & \SFiveEssMeanAIDramatic & \SFiveEssMeanUnderOneHour & \SFiveEssMeanTenHours & \SFiveEssMeanHundredHours\\
Creativity & \SFiveCreatMeanHumanNone & \SFiveCreatMeanAINone & \SFiveCreatMeanHumanSlight & \SFiveCreatMeanAISlight & \SFiveCreatMeanHumanDramatic & \SFiveCreatMeanAIDramatic & \SFiveCreatMeanUnderOneHour & \SFiveCreatMeanTenHours & \SFiveCreatMeanHundredHours\\
\bottomrule
\end{tabular}

\smallskip
\begin{minipage}{\textwidth}
\footnotesize\raggedright
Estimates come from the registered linear mixed model, which gives each participant and each original its own baseline. $df$: degrees of freedom by Satterthwaite's approximation; the rescaled tests have the same numerator degrees of freedom. The rescaled analysis includes the two six-point transformativeness items, which the primary models leave out; essence change and creativity have no six-point items, so their rescaled tests repeat the primary ones on a 0-to-1 scale. An equivalence test passes when the 90\% confidence interval (CI) lies wholly inside the bound, which means small, not zero. Human, AI: the label of the original. Supplementary Table~\ref{tab:supp-labeleffects} gives the label's effect at each level with its registered verdict.
\end{minipage}
\end{table}


\paragraph*{The label's penalty at each editing time, on every measure.}
Supplementary Note~5 explains the exploratory breakdown of the creativity interaction. Supplementary Table~\ref{tab:suppB-s5breakdown} extends it to transformativeness and essence change, whose registered interactions were not significant. On both, as on creativity, the largest penalties at Slight and Dramatic fell in the 10-hour groups. Because editing time varied between participants, each column of that table compares separate groups, and one unusual 10-hour group could produce this pattern on all three measures.

% Study 5, exploratory: the label's penalty at each stated editing time and level, on all three
% measures, from the cross-study refit of Study 5's registered model.
% Every number is a macro from generated/numbers.tex; supplement/tables/README_suppB.md lists
% the pipeline tables behind them.
\begin{table}[htbp]
\centering
\footnotesize
\caption{\textbf{Exploratory: Study~5's label penalty at each stated editing time and level.} Each entry is the predicted mean rating under the Human label minus the predicted mean under the AI label, in points on the 1-to-5 scale, with its two-sided $p$-value, not corrected for the 27 comparisons. A positive value means that the AI label lowered the rating. Editing time varied between participants, so each column compares two separate groups of \SFiveNCellMin{} to \SFiveNCellMax{} people. The breakdown was specified on 9 October 2026, after the registered results were known, and it changes no registered verdict. Among the registered tests behind it, only the creativity interaction was significant (Supplementary Table~\ref{tab:suppB-s5}a).}\label{tab:suppB-s5breakdown}
\setlength{\tabcolsep}{6pt}
\begin{tabular}{@{}lccc@{}}
\toprule
 & \multicolumn{3}{c}{Stated editing time}\\
\cmidrule(l){2-4}
Level & Under 1 hour & 10 hours & 100 hours\\
\midrule
\multicolumn{4}{@{}l}{\textit{Transformativeness}}\\
None & \XSTimeCheckSFiveTransPenaltyUnderOneHourNoneEst{} ($p$~\XSTimeCheckSFiveTransPenaltyUnderOneHourNoneP) & \XSTimeCheckSFiveTransPenaltyTenHoursNoneEst{} ($p$~\XSTimeCheckSFiveTransPenaltyTenHoursNoneP) & \XSTimeCheckSFiveTransPenaltyHundredHoursNoneEst{} ($p$~\XSTimeCheckSFiveTransPenaltyHundredHoursNoneP)\\
Slight & \XSTimeCheckSFiveTransPenaltyUnderOneHourSlightEst{} ($p$~\XSTimeCheckSFiveTransPenaltyUnderOneHourSlightP) & \XSTimeCheckSFiveTransPenaltyTenHoursSlightEst{} ($p$~\XSTimeCheckSFiveTransPenaltyTenHoursSlightP) & \XSTimeCheckSFiveTransPenaltyHundredHoursSlightEst{} ($p$~\XSTimeCheckSFiveTransPenaltyHundredHoursSlightP)\\
Dramatic & \XSTimeCheckSFiveTransPenaltyUnderOneHourDramaticEst{} ($p$~\XSTimeCheckSFiveTransPenaltyUnderOneHourDramaticP) & \XSTimeCheckSFiveTransPenaltyTenHoursDramaticEst{} ($p$~\XSTimeCheckSFiveTransPenaltyTenHoursDramaticP) & \XSTimeCheckSFiveTransPenaltyHundredHoursDramaticEst{} ($p$~\XSTimeCheckSFiveTransPenaltyHundredHoursDramaticP)\\
\addlinespace[3pt]
\multicolumn{4}{@{}l}{\textit{Essence change}}\\
None & \XSTimeCheckSFiveEssPenaltyUnderOneHourNoneEst{} ($p$~\XSTimeCheckSFiveEssPenaltyUnderOneHourNoneP) & \XSTimeCheckSFiveEssPenaltyTenHoursNoneEst{} ($p$~\XSTimeCheckSFiveEssPenaltyTenHoursNoneP) & \XSTimeCheckSFiveEssPenaltyHundredHoursNoneEst{} ($p$~\XSTimeCheckSFiveEssPenaltyHundredHoursNoneP)\\
Slight & \XSTimeCheckSFiveEssPenaltyUnderOneHourSlightEst{} ($p$~\XSTimeCheckSFiveEssPenaltyUnderOneHourSlightP) & \XSTimeCheckSFiveEssPenaltyTenHoursSlightEst{} ($p$~\XSTimeCheckSFiveEssPenaltyTenHoursSlightP) & \XSTimeCheckSFiveEssPenaltyHundredHoursSlightEst{} ($p$~\XSTimeCheckSFiveEssPenaltyHundredHoursSlightP)\\
Dramatic & \XSTimeCheckSFiveEssPenaltyUnderOneHourDramaticEst{} ($p$~\XSTimeCheckSFiveEssPenaltyUnderOneHourDramaticP) & \XSTimeCheckSFiveEssPenaltyTenHoursDramaticEst{} ($p$~\XSTimeCheckSFiveEssPenaltyTenHoursDramaticP) & \XSTimeCheckSFiveEssPenaltyHundredHoursDramaticEst{} ($p$~\XSTimeCheckSFiveEssPenaltyHundredHoursDramaticP)\\
\addlinespace[3pt]
\multicolumn{4}{@{}l}{\textit{Creativity}}\\
None & \XSTimeCheckSFiveCreatPenaltyUnderOneHourNoneEst{} ($p$~\XSTimeCheckSFiveCreatPenaltyUnderOneHourNoneP) & \XSTimeCheckSFiveCreatPenaltyTenHoursNoneEst{} ($p$~\XSTimeCheckSFiveCreatPenaltyTenHoursNoneP) & \XSTimeCheckSFiveCreatPenaltyHundredHoursNoneEst{} ($p$~\XSTimeCheckSFiveCreatPenaltyHundredHoursNoneP)\\
Slight & \XSTimeCheckSFiveCreatPenaltyUnderOneHourSlightEst{} ($p$~\XSTimeCheckSFiveCreatPenaltyUnderOneHourSlightP) & \XSTimeCheckSFiveCreatPenaltyTenHoursSlightEst{} ($p$~\XSTimeCheckSFiveCreatPenaltyTenHoursSlightP) & \XSTimeCheckSFiveCreatPenaltyHundredHoursSlightEst{} ($p$~\XSTimeCheckSFiveCreatPenaltyHundredHoursSlightP)\\
Dramatic & \XSTimeCheckSFiveCreatPenaltyUnderOneHourDramaticEst{} ($p$~\XSTimeCheckSFiveCreatPenaltyUnderOneHourDramaticP) & \XSTimeCheckSFiveCreatPenaltyTenHoursDramaticEst{} ($p$~\XSTimeCheckSFiveCreatPenaltyTenHoursDramaticP) & \XSTimeCheckSFiveCreatPenaltyHundredHoursDramaticEst{} ($p$~\XSTimeCheckSFiveCreatPenaltyHundredHoursDramaticP)\\
\bottomrule
\end{tabular}

\smallskip
\begin{minipage}{\textwidth}
\footnotesize\raggedright
The estimates come from a refit of Study~5's registered model in the cross-study analysis, which checks that, averaged over editing time with equal weights, the refit reproduces Study~5's registered penalty at each level. None is the original shown again unchanged, where ratings sit near the floor of the scale. Supplementary Figure~\ref{figS:editingtime} plots the creativity rows.
\end{minipage}
\end{table}


\paragraph*{Other terms of the registered model.}
The registered model contains terms that the registration names no test of. Examined as exploratory, uncorrected tests, the term that lets the label's effect differ by level was significant on transformativeness and essence change, the label's overall effect was significant on creativity, and the term that lets the effect of editing time differ by level fell just under .05 on creativity. The registration's own test of the label is the set of equivalence tests at each level in Supplementary Table~\ref{tab:supp-labeleffects}, not this overall term. Supplementary Note~5 reports the three-way term on creativity. Their statistics, apart from the three-way term, are in the Study~5 tables that the analysis code writes (main text, Code availability).

\paragraph*{Participants who wrote that a machine made the work.}
A fixed word list, run without access to labels or ratings, flagged explanations stating that a machine made the work (exploratory count). Of the \SFiveMachineMadeNHuman{} Human-label participants, \SFiveMachineMadeN{} wrote this at least once, in \SFiveMachineMadeResponses{} of their \SFiveMachineMadeNResponses{} explanations. This is a count of statements, not of disbelief, and Study~5 registered no refit without these participants.

\paragraph*{Registered text checks.}
Two checks on the creativity explanations were registered for every study, and Study~5 ran both (registered secondary; Section~\ref{sec:suppA-text} describes both). In the first, ridge regression, a regression that shrinks each word's weight toward zero, predicted each rating and each participant's label from the words of explanations written by participants left out of the fit. The words predicted the rating: predicted and actual ratings correlated at \SFiveTextRecRatingCreat{} (95\% confidence interval, CI, the range of values the data do not rule out at the 95\% level: \SFiveTextRecRatingCreatLo{} to \SFiveTextRecRatingCreatHi), and in runs with the ratings shuffled, 95\% of correlations stayed below \SFiveTextRecRatingCreatNull. The words showed the label, barely. The AUC, the chance that a randomly chosen AI-label explanation scores above a randomly chosen Human-label one, so that 0.5 is chance, was \SFiveTextRecArmCreat{} (\SFiveTextRecArmCreatLo{} to \SFiveTextRecArmCreatHi), with a lower end just above chance before rounding; the rating alone gave \SFiveTextRecArmFromRatingCreat. Adding the words to the rating changed the AUC by \SFiveTextRecIncrCreat{} (\SFiveTextRecIncrCreatLo{} to \SFiveTextRecIncrCreatHi), which is inconclusive. The second check predicted that AI-label participants would reuse their own words across explanations more than Human-label participants, measured as a person's word overlap among their own explanations minus their overlap with other people's explanations of the same second work, with explanation length held constant. The difference, AI minus Human, was \SFiveTextOverlapCreatDiff{} (\SFiveTextOverlapCreatLo{} to \SFiveTextOverlapCreatHi; Holm $p$~\SFiveTextOverlapCreatPHolm): in the predicted direction but inconclusive, which is not evidence that the groups reuse words equally. Every explanation was written after its rating, so these checks show what the words track, not why anyone rated as they did.

\paragraph*{Exploratory text analyses.}
The registration names the remaining text analyses as exploratory. Word tone and concreteness, scored from published norms, did not differ by label after a Benjamini--Hochberg correction, which limits the expected share of false positives among results called significant. No topic found without labels, a group of words that tend to occur together, differed in weight by label after the same correction. Before correction, AI-label explanations leaned toward topics about the second work looking identical or exact, and Human-label explanations toward a topic about changes in essence. In predicting the rating, words for the size of a change, such as `dramatically' and `complete', and words of praise went with higher ratings, and `identical', `exact' and `copy' with lower ones. Their estimates are in the Study~5 tables that the analysis code writes (main text, Code availability).

\clearpage
